% ============================================================
%  附录G 薛定谔方程的变分原理
%  底本：高等教育出版社《固体理论》第二版；原书印刷页 557–558
%  （PDF 页 = 印刷页 + 20）
%  公式编号自检：(G.1)–(G.15)，自动编号 15 个，无缺号
%  （任务单所称缺 [11] 系 OCR 把 (G.11) 的 tag 残缺为"(.11)"
%  未被脚本计入，原书 p558 在号）。
%
%  已修正的 OCR/原书错误（除全角句点转"."、收紧 \mathrm{d}、
%  拆开行内公式与中文等常规项外）：
%  1. [定理]句及"将 φ(r) 代入能量泛函"OCR 作 φ(r)，按 p557
%     原书复原为 ψ(r)。
%  2. (G.3) 势能项与分母 OCR 作 |φ(r)|²、|φ,(r)|²，按 p557
%     原书复原为 |ψ(r)|²（原书全式均排 ψ）。
%  3. (G.11) 前一未编号式 OCR 的"∇φ_k、φ_k"按 p557–558 复原为
%     ∇ψ_k、ψ_k；该块首行"F[∮φ,ψ_k]"乱码按 p557 复原为
%     F[φ,ψ_k]。
%  4. (G.11) tag 残缺"(.11)"已归还本式；其首行括号
%     "( −∇²u_k−2ik·∇u_k+k²u_k )"OCR 乱码按 p558 复原。
%  5. (G.13)(G.14) OCR 作 φ、φ_k、δ̂φ，按 p558 原书复原为
%     ψ、ψ_k、δψ（原书两式均排 ψ，δ 无帽）。
%  6. (G.4) 原书印"0(δψ)²"（数字零），(G.13)–(G.15) 印
%    "O(δψ)²"（字母 O），照原书分别照录，未统一。
%  正文引用 (1.5.13)、(1.5.22) 照录（指第一章式号）。
%  遗留问题：无
% ============================================================
\chapter{薛定谔方程的变分原理}

[定理] 设 $\psi(\boldsymbol{r})$ 是满足布洛赫定理的一次可微变分函数
\begin{equation}
  \psi=\psi_{k}+\delta\psi
\end{equation}
$\delta\psi$ 为小量，$\psi_{k}$ 是周期场中单电子薛定谔方程的解
\begin{equation}
  -\frac{\hbar^{2}}{2m}\nabla^{2}\psi_{k}
  +V(\boldsymbol{r})\psi_{k}=E_{k}\psi_{k}
\end{equation}
代表 $k$ 波矢和 $E_{k}$ 能量的布洛赫函数.将 $\psi(\boldsymbol{r})$ 代入能量泛函
\begin{equation}
  E[\psi]=\frac{\displaystyle\int\mathrm{d}\tau\left(
  \frac{\hbar^{2}}{2m}\left|\nabla\psi(\boldsymbol{r})\right|^{2}
  +V(\boldsymbol{r})\left|\psi(\boldsymbol{r})\right|^{2}\right)}
  {\displaystyle\int\mathrm{d}\tau
  \left|\psi(\boldsymbol{r})\right|^{2}}
\end{equation}
可证明
\begin{equation}
  E[\psi]=E[\psi_{k}]+0\,(\delta\psi)^{2}
\end{equation}
其中
\begin{equation}
  E[\psi_{k}]=E_{k}
\end{equation}
由于(G.4)中不含 $\delta\psi$ 一次项，说明能量泛函 $E[\psi]$ 对于式(G.1)所规定的 $\psi$ 是稳定的，而式(G.5)表明，稳定点的能量泛函值就是薛定谔方程(G.2)的能量本征值.

[证明] 为使记号简单起见定义
\begin{equation}
  F[\phi,\chi]\equiv\int\mathrm{d}\tau\left(
  \frac{\hbar^{2}}{2m}\nabla\phi^{*}\cdot\nabla\chi
  +V(\boldsymbol{r})\phi^{*}\chi\right)
\end{equation}
\begin{equation}
  (\phi,\chi)\equiv\int\mathrm{d}\tau\,\phi^{*}\chi
\end{equation}
这样 $E[\psi]$ 简记为
\begin{equation}
  E[\psi]=\frac{F[\psi,\psi]}{(\psi,\psi)}
\end{equation}
设 $\phi$ 为满足 $k$ 波矢的布洛赫定理的任意函数
\begin{equation}
  \phi(\boldsymbol{r})=v(\boldsymbol{r})
  \mathrm{e}^{\mathrm{i}\boldsymbol{k}\cdot\boldsymbol{r}}
\end{equation}
而
\begin{equation}
  \psi_{k}(\boldsymbol{r})=u_{k}(\boldsymbol{r})
  \mathrm{e}^{\mathrm{i}\boldsymbol{k}\cdot\boldsymbol{r}}
\end{equation}
其中 $u_{k}(\boldsymbol{r})$ 与 $v(\boldsymbol{r})$ 都是正点阵的周期函数，于是
\begin{equation*}
\begin{aligned}
  F[\phi,\psi_{k}]
  &=\int\mathrm{d}\tau\left(\frac{\hbar^{2}}{2m}
  \nabla\phi^{*}\cdot\nabla\psi_{k}
  +V\phi^{*}\psi_{k}\right)\\
  &=\frac{\hbar^{2}}{2m}\Bigl\{k^{2}\int\mathrm{d}\tau\,
  v^{*}u_{k}
  +\mathrm{i}\boldsymbol{k}\cdot\int\mathrm{d}\tau\,
  \left(u_{k}\nabla v^{*}-v^{*}\nabla u_{k}\right)\\
  &\qquad+\int\mathrm{d}\tau\,\nabla v^{*}\cdot\nabla u_{k}
  \Bigr\}+\int\mathrm{d}\tau\,V\phi^{*}\psi_{k}
\end{aligned}
\end{equation*}
对于周期函数 $v^{*}$ 和 $u_{k}$ 利用式(1.5.13)，对于 $\nabla v^{*}$ 和 $\nabla u_{k}$ 应用式(1.5.22)，求得
\begin{equation}
\begin{aligned}
  F[\phi,\psi_{k}]
  &=\frac{\hbar^{2}}{2m}\int\mathrm{d}\tau\,v^{*}
  \left(-\nabla^{2}u_{k}
  -2\mathrm{i}\boldsymbol{k}\cdot\nabla u_{k}
  +k^{2}u_{k}\right)
  +\int\mathrm{d}\tau\,V\phi^{*}\psi_{k}\\
  &=-\frac{\hbar^{2}}{2m}\int\mathrm{d}\tau\,\phi^{*}
  \nabla^{2}\psi_{k}
  +\int\mathrm{d}\tau\,V\phi^{*}\psi_{k}\\
  &=\int\mathrm{d}\tau\,\phi^{*}\left[-\frac{\hbar^{2}}{2m}
  \nabla^{2}+V(\boldsymbol{r})\right]\psi_{k}\\
  &=E_{k}\left(\phi,\psi_{k}\right)
\end{aligned}
\end{equation}
类似的方法可证明
\begin{equation}
  F[\psi_{k},\phi]=E_{k}\left(\psi_{k},\phi\right)
\end{equation}
现将式(G.1)代入式(G.6)，再利用式(G.11)与式(G.12)，求得
\begin{equation}
\begin{aligned}
  F[\psi,\psi]
  &=F[\psi_{k}+\delta\psi,\psi_{k}+\delta\psi]\\
  &=F[\psi_{k},\psi_{k}]+F[\delta\psi,\psi_{k}]
  +F[\psi_{k},\delta\psi]+O\left(\delta\psi\right)^{2}\\
  &=E_{k}\left\{(\psi_{k},\psi_{k})+(\psi_{k},\delta\psi)
  +(\delta\psi,\psi_{k})\right\}
  +O\left(\delta\psi\right)^{2}
\end{aligned}
\end{equation}
又由于
\begin{equation}
  (\psi,\psi)=(\psi_{k},\psi_{k})+(\psi_{k},\delta\psi)
  +(\delta\psi,\psi_{k})+O\left(\delta\psi\right)^{2}
\end{equation}
最后，将式(G.13)与式(G.14)代入式(G.8)并作小量展开
\begin{equation}
  E[\psi]=\frac{F[\psi,\psi]}{(\psi,\psi)}
  =E_{k}+O\left(\delta\psi\right)^{2}
\end{equation}
当式(G.15)中取 $\delta\psi=0$，即得到 $E[\psi_{k}]=E_{k}$，于是变分原理式(G.4)与(G.5)得证.

应当特别指出，在证明过程并未要求变分函数 $\psi$ 的一阶导数连续（二次可微条件），因此本节的原理适用于对缀加平面波作变分计算.
